% TOP_PROBLEM: {"id": 30002459, "title": "Sidorenko's Inequality for Bipartite Graphs", "category_id": 3, "category": {"id": 3, "name": "graph_theory", "display_name": "Graph Theory", "description": "Problems involving graphs, networks, and their properties.", "slug": "graph-theory", "order_index": 3, "created_at": "2026-07-31T15:26:25.671Z"}, "status": "open"}
% FIRST_POSTED: 2026-10-06
% ATTEMPT_STATUS: solved
% OPENAI_CLAIM: {"schema_version":1,"source_repo":"https://github.com/openai/math","source_commit":"adc7f1241b42e322a6451854ab7e4b4c146bf78a","release_date":"2026-10-06","match":"full","resolution":"disproved","independently_reviewed":false,"second_pass":"agreed","families":[{"family":"161","title":"Counterexamples to Sidorenko’s conjecture and the forcing conjecture","subject":"Combinatorics","manuscripts":[{"title":"A counterexample to Sidorenko's conjecture","pdf_path":"preprints/A-counterexample-to-Sidorenkos-conjecture-September-23-2026/paper.pdf","readme_path":"preprints/A-counterexample-to-Sidorenkos-conjecture-September-23-2026/README.md","date":"2026-09-23","theorem_ref":"thm:main (sections/introduction.tex:18); cor:amplification (sections/introduction.tex:279); cor:forcing (sections/introduction.tex:335)","statement_sources":[{"path":"preprints/A-counterexample-to-Sidorenkos-conjecture-September-23-2026/build/sections/introduction.tex","line":18,"label":"thm:main"},{"path":"preprints/A-counterexample-to-Sidorenkos-conjecture-September-23-2026/build/sections/introduction.tex","line":279,"label":"cor:amplification"},{"path":"preprints/A-counterexample-to-Sidorenkos-conjecture-September-23-2026/build/sections/introduction.tex","line":335,"label":"cor:forcing"}]}],"lean":{"doc_path":"lean/docs/161.md","comparators":["lean/ComparatorChallenges/SidorenkoCounterexample.json","lean/ComparatorChallenges/SidorenkoCounterexample.lean"],"declarations":[],"covers_main_theorem":true},"reasoning_trace":null}]}
\documentclass[11pt]{article}
\usepackage[margin=25mm]{geometry}
\usepackage{fontspec}
\setmainfont{Latin Modern Roman}
\usepackage{amsmath,amssymb,mathtools}
\usepackage{unicode-math}
\setmathfont{Latin Modern Math}
\usepackage{enumitem}
\usepackage{xurl,hyperref,bookmark}
\hypersetup{colorlinks=true,urlcolor=blue}
\setcounter{secnumdepth}{0}
\setlength{\parindent}{0pt}
\setlength{\parskip}{5pt}
\setlength{\emergencystretch}{3em}
\begin{document}
\section*{Sidorenko's Inequality for Bipartite Graphs}

\subsection{OpenAI's claim}
\textit{Sidorenko and forcing counterexamples}: A connected bipartite graph with 35 vertices and 66 edges is claimed to violate Sidorenko’s homomorphism-density lower bound in a finite host. The same pattern is also claimed not to force quasirandomness when the edge and pattern densities match their random values.

\subsection{Mathematical statement of the claimed result}
\paragraph{Mathematical claim 1.} Let $H$ be the incidence graph of the explicitly tabulated $22$ triples on $13$ points: it is connected and bipartite with $35$ vertices and $66$ edges. For $t(H,G)=\operatorname{hom}(H,G)/|V(G)|^{35}$ and $p(G)=2|E(G)|/|V(G)|^2$, a finite simple graph $G$ with $p(G)>0$ is claimed to satisfy $t(H,G)<p(G)^{66}$. Moreover, for every $c>0$ another graph $F$ has $0<t(K_2,F)$ and $t(H,F)<c\,p(F)^{66}$. The forcing corollary claims some $p\in(0,1)$ and a nonconstant symmetric measurable graphon $W:[0,1]^2\to[0,1]$ with \[t(K_2,W)=p,\qquad t(H,W)=p^{66},\qquad t(H,W)=\int_{[0,1]^{35}}\prod_{ij\in E(H)}W(x_i,x_j)\,dx.\] Thus equality at that density does not force this cyclic bipartite graph to be quasirandom.

\subsection{Outline of the argument}
\textit{Sidorenko and forcing counterexamples}: The construction combines rank and discriminant relations for symmetric matrices with counts of symplectic intersections. Quadratic Gauss-sum estimates control joint signs, while type scores distinguish the desired vertex roles. Singular configurations are estimated separately so that the signed rank-layer correlations survive normalization and transfer to a finite host graph.

\subsection{Scope relative to this problem}
\textit{Sidorenko and forcing counterexamples}: The explicit finite bipartite graph H and host G satisfy t(H,G)<t(K\_2,G)\textasciicircum{}e(H). Passing to the associated step graphon also refutes the graphon formulation.

The claimed counterexample lies within the question’s hypotheses and contradicts its universal conclusion.

\subsection{Formal verification}
Family 161 describes Lean coverage for the fixed-pattern Sidorenko deficit; the forcing consequence is not selected.

\textit{Sidorenko and forcing counterexamples}: The manuscript links \texttt{SidorenkoCounterexample.lean}.

The documented selection includes this manuscript’s main assertion; that assertion retains the restrictions stated above.

This site has not run the Lean development. These links describe OpenAI’s selected formal statements and should be read with their stated scope.

\subsection{Status caveat}
OpenAI reports that an internal model produced the manuscripts. Their arguments have not been independently verified as proofs by this site. Any portion outside the documented formal scope may contain mistakes, and the catalogue label reports the announced mathematical claim.

\subsection{Sources}
\begin{itemize}
\item OpenAI, \emph{A counterexample to Sidorenko's conjecture}, 2026-09-23; citation key \texttt{OAI:A-counterexample-to-Sidorenkos-conjecture-September-23-2026}. \href{https://github.com/openai/math/blob/adc7f1241b42e322a6451854ab7e4b4c146bf78a/preprints/A-counterexample-to-Sidorenkos-conjecture-September-23-2026/paper.pdf}{PDF (GitHub)}; \href{https://github.com/openai/math/raw/adc7f1241b42e322a6451854ab7e4b4c146bf78a/preprints/A-counterexample-to-Sidorenkos-conjecture-September-23-2026/paper.pdf}{PDF (direct)}; \href{https://github.com/openai/math/blob/adc7f1241b42e322a6451854ab7e4b4c146bf78a/preprints/A-counterexample-to-Sidorenkos-conjecture-September-23-2026/README.md}{Manuscript page and citation}; \href{https://github.com/openai/math/blob/main/preprints/A-counterexample-to-Sidorenkos-conjecture-September-23-2026/paper.pdf}{Latest version}; Statement source: \href{https://github.com/openai/math/blob/adc7f1241b42e322a6451854ab7e4b4c146bf78a/preprints/A-counterexample-to-Sidorenkos-conjecture-September-23-2026/build/sections/introduction.tex\#L18}{thm:main}; Statement source: \href{https://github.com/openai/math/blob/adc7f1241b42e322a6451854ab7e4b4c146bf78a/preprints/A-counterexample-to-Sidorenkos-conjecture-September-23-2026/build/sections/introduction.tex\#L279}{cor:amplification}; Statement source: \href{https://github.com/openai/math/blob/adc7f1241b42e322a6451854ab7e4b4c146bf78a/preprints/A-counterexample-to-Sidorenkos-conjecture-September-23-2026/build/sections/introduction.tex\#L335}{cor:forcing}.
\item Family 161: \href{https://github.com/openai/math/blob/adc7f1241b42e322a6451854ab7e4b4c146bf78a/lean/docs/161.md}{Lean scope}.
\item \href{https://github.com/openai/math/blob/adc7f1241b42e322a6451854ab7e4b4c146bf78a/lean/ComparatorChallenges/SidorenkoCounterexample.json}{Comparator SidorenkoCounterexample.json}.
\item \href{https://github.com/openai/math/blob/adc7f1241b42e322a6451854ab7e4b4c146bf78a/lean/ComparatorChallenges/SidorenkoCounterexample.lean}{Comparator SidorenkoCounterexample.lean}.
\item \href{https://github.com/openai/math/blob/adc7f1241b42e322a6451854ab7e4b4c146bf78a/CONTENTS.md}{OpenAI catalogue entry}.
\end{itemize}
\end{document}
